% ================================================================
% Cover Letter — International Journal of Fatigue
% ================================================================
\documentclass[12pt,a4paper]{article}
\usepackage[margin=2.5cm]{geometry}
\usepackage{hyperref}

\begin{document}

\noindent Dear Editor,

\vspace{6pt}

\noindent I am pleased to submit our manuscript entitled
``Surface Roughness and Standard Choice Drive Gear Bending Fatigue
Life Prediction Divergence: A Full-Factorial Variance Decomposition''
for consideration for publication in the \textit{International Journal
of Fatigue}.

\vspace{6pt}

\noindent\textbf{What this paper does.}
We apply a full-factorial ANOVA framework to decompose the total
variance in gear tooth bending fatigue life prediction into six
methodological sources. Holding geometry, material, and loading
constant, we systematically vary the choices an engineer must make:
which standard (ISO~6336, AGMA~2001, FKM), which surface roughness
grade (Rz~=~4, 15, 50~$\mu$m), which S--N curve (MatWeb, MMPDS), which
mean-stress correction (Goodman, Gerber, Morrow, SWT), and so on. The
analysis reveals that mean-stress correction method (39.0\%) and
standard choice (34.5\%) together account for nearly three-quarters of
all prediction divergence---a finding that contradicts the common
assumption that mean-stress effects are minor for gear bending at
$R=0$. All formulas in the fatigue model have been verified against the
original standards (ISO~6336-3:2019, AGMA~2001-D04, FKM~7th~ed.).

\vspace{6pt}

\noindent\textbf{Why IJ Fatigue.}
The paper addresses a core question in fatigue life prediction---how
much does methodological choice matter relative to physical
uncertainty?---using a quantitative variance-decomposition framework
that has been applied in other fields but is new to gear fatigue. The
analysis is purely computational (no experimental data), uses only
open-source tools and publicly available material databases, and is
fully reproducible. We believe the methodology and findings are
directly relevant to the journal's readership in fatigue life
prediction, standards development, and uncertainty quantification.

\vspace{6pt}

\noindent\textbf{Novelty and significance.}
To our knowledge, this is the first study to deploy a controlled
full-factorial ANOVA to rank the sources of method-induced divergence
in gear bending fatigue life prediction. Prior work has compared
individual factors in isolation (e.g., Goodman vs.\ Gerber, or ISO vs.\
AGMA) but has not systematically mapped the full interaction structure.
The finding that mean-stress correction is the largest source of
variance at $R=0$---contradicting engineering intuition---is both
methodologically novel and practically significant.

\vspace{6pt}

\noindent The manuscript has not been published and is not under
consideration elsewhere. All authors have approved the submission.

\vspace{6pt}

\noindent\textbf{AI assistance disclosure.}
Language editing and code development for this work were assisted by
large language models (Claude, Anthropic; Hermes Agent). All numerical
results were generated by executing the open-source Python pipeline
provided with the submission. Every formula in the fatigue model has
been verified against the original standards (ISO~6336-3:2019,
AGMA~2001-D04, FKM~7th~ed.). The author takes full responsibility for
the scientific content.

\vspace{12pt}

\noindent Sincerely,

\vspace{18pt}

\noindent Zhilei Chen \\
Guangdong Peizheng College \\
Data and Computer Science \\
53 Peizheng Rd., Chini Town, Huadu \\
Guangzhou, Guangdong 510830, China \\
E-mail: 2604513@peizheng.edu.cn

\end{document}
